\begin{pastexam}[label={px:ntu-115-dsa-18}]{臺大資工所 115 學年度 第 18 題}
（單選）

(4 points) Which of the following properties is true for \textit{some} optimal prefix-free binary code and directly justifies the greedy choice made in the Huffman algorithm?
\begin{enumerate}[label=(\Alph*)]
\item The two symbols with the lowest frequencies must appear at the same depth in the code tree.
\item The symbol with the lowest frequency must be assigned the longest codeword, but the symbol with the second smallest frequency need not be.
\item All leaves at the maximum depth must correspond to symbols with the lowest frequency.
\item The two symbols with the lowest frequencies must have equal codeword lengths, but they need not share the same parent.
\item The two symbols with the lowest frequencies must be siblings at the maximum depth of the code tree.
\end{enumerate}
\end{pastexam}

\begin{pastsoln}
答案是 \textbf{(E)}。

霍夫曼演算法第一步把頻率最小的兩個符號 $x$、$y$ 接到同一個新節點底下，也就是讓它們成為兄
弟。 要說明這一步不會錯，需要的正是「存在一棵最佳碼樹，$x$、$y$ 是兄弟，而且都在最深的一
層」：這就是資結書引理~19.3。 有了它，把 $x$、$y$ 收成一個符號之後，剩下的問
題的最佳解展開回去就是原問題的最佳解（資結書引理~19.4）。 逐項看：

\begin{itemize}
\item (A)：對某一棵最佳碼樹成立（(E) 成立時 (A) 也成立），但「一樣深」不代表是兄弟。 霍夫曼
把 $x$、$y$ 合成同一個節點，需要的是兄弟關係，所以 (A) 不足以支持這個選擇。
\item (B)：前半句對某棵最佳碼樹成立，但它只講頻率最小的那一個；後半句還說第二小的不必在最深
層。 霍夫曼同時處理最小的兩個，(B) 支持不了。
\item (C)：錯，連「某一棵最佳碼樹」都找不到。 例如頻率 $1, 2, 2, 2$：$4$ 個葉節點的嚴格二元
樹只有兩種形狀，碼長 $(2, 2, 2, 2)$ 或 $(1, 2, 3, 3)$ 的某種排列。 前者總位元數
$(1 + 2 + 2 + 2) \cdot 2 = 14$；後者最好是頻率 $1$ 和一個 $2$ 放在深度 $3$，
$1 \cdot 3 + 2 \cdot 3 + 2 \cdot 2 + 2 \cdot 1 = 15$。 所以唯一的最佳形狀是四個碼長都是 $2$，
最深的一層有頻率 $2$ 的符號，它們不是頻率最小的。
\item (D)：碼長相等對某棵最佳碼樹成立，但它明說「不必有同一個父節點」，正好不是霍夫曼需要的
性質。
\item (E)：正確，就是資結書引理~19.3，證明用交換論證：任取一棵最佳碼樹，把
$x$、$y$ 和最深的一對兄弟葉節點交換，總位元數不會變大。
\end{itemize}
\end{pastsoln}
